% !TEX program = xelatex
% Circuitikz 示例：IC 设计常用电路元件
% 编译方式：XeLaTeX

\documentclass[a4paper, 12pt]{article}
\usepackage{fontspec}
\usepackage{xeCJK}
\setCJKmainfont[BoldFont={STSong},BoldFeatures={FakeBold=2}]{STSong}
\usepackage[margin=2cm]{geometry}
\usepackage{circuitikz}
\usepackage{subcaption}
\usepackage{xcolor}

% 颜色定义
\definecolor{icblue}{RGB}{0, 70, 140}

\begin{document}

\section*{Circuitikz 元件速查（IC 设计常用）}

% ============================================================
% 示例1：基本 MOS 管符号
% ============================================================
\subsection*{1. 基本 MOS 管}

\begin{circuitikz}[american, scale=1.2, transform shape]
  % NMOS
  \draw (0,0) node[nmos] (n1) {NMOS};
  \draw (n1.gate) -- ++(-1,0) node[left]{$V_{in}$};
  \draw (n1.drain) -- ++(0,0.5) node[vcc]{$V_{DD}$};
  \draw (n1.source) -- ++(0,-0.5) node[ground]{};

  % PMOS
  \draw (4,0) node[pmos] (p1) {PMOS};
  \draw (p1.gate) -- ++(-1,0) node[left]{$V_{in}$};
  \draw (p1.drain) -- ++(0,-0.5) node[ground]{};
  \draw (p1.source) -- ++(0,0.5) node[vcc]{$V_{DD}$};
\end{circuitikz}

\vspace{1cm}

% ============================================================
% 示例2：电流镜
% ============================================================
\subsection*{2. NMOS 电流镜}

\begin{circuitikz}[american, scale=1.0, transform shape]
  % 参考管 M1
  \draw (0,0) node[nmos] (m1) {$M_1$};
  \draw (m1.gate) -- (m1.drain) -- ++(0,0.3) coordinate (gate);
  \draw (m1.source) -- ++(0,-0.3) node[ground]{};

  % 输出管 M2
  \draw (2.5,0) node[nmos] (m2) {$M_2$};
  \draw (m2.gate) -- ++(0,0.3) -- (gate) -- (m1.gate);
  \draw (m2.source) -- ++(0,-0.3) node[ground]{};

  % 电流源
  \draw (m1.drain) -- ++(0,0.8) to[I, l=$I_{ref}$] ++(0,1.2) node[vcc]{$V_{DD}$};
  \draw (m2.drain) -- ++(0,0.8) -- ++(0.5,0) node[right]{$I_{out}$};
  \draw (m2.drain) -- ++(0,2) node[vcc]{$V_{DD}$};
\end{circuitikz}

\vspace{1cm}

% ============================================================
% 示例3：差分对（运放输入级）
% ============================================================
\subsection*{3. NMOS 差分对（运放输入级）}

\begin{circuitikz}[american, scale=1.0, transform shape]
  % M1
  \draw (0,0) node[nmos] (m1) {$M_1$};
  \draw (m1.gate) -- ++(-1.2,0) node[left]{$V_{in+}$};
  \draw (m1.drain) -- ++(0,0.5) coordinate (d1);

  % M2
  \draw (2.5,0) node[nmos] (m2) {$M_2$};
  \draw (m2.gate) -- ++(1.2,0) node[right]{$V_{in-}$};
  \draw (m2.drain) -- ++(0,0.5) coordinate (d2);

  % 共源节点
  \draw (m1.source) -- ++(0,-0.3) coordinate (ss);
  \draw (m2.source) -- (ss);

  % 尾电流源
  \draw (ss) -- ++(0,-0.3) node[nmos, anchor=gate, rotate=90] (mb) {$M_b$};
  \draw (mb.drain) -- (ss);
  \draw (mb.source) -- ++(0,-0.3) node[ground]{};
  \draw (mb.gate) -- ++(-0.8,0) node[left]{$V_{bias}$};

  % 负载（有源负载简化）
  \draw (d1) -- ++(0,1.0) to[R, l=$R_D$] ++(0,1.2) node[vcc]{$V_{DD}$};
  \draw (d2) -- ++(0,1.0) to[R, l=$R_D$] ++(0,1.2) node[vcc]{$V_{DD}$};

  % 输出
  \draw (d2) -- ++(1,0) node[right]{$V_{out}$};
\end{circuitikz}

\vspace{1cm}

% ============================================================
% 示例4：简单 LDO 结构框图
% ============================================================
\subsection*{4. LDO 稳压器基本结构}

\begin{circuitikz}[american, scale=0.9, transform shape]
  % 误差放大器
  \draw (-2,0) node[op amp] (amp) {};
  \draw (amp.-) -- ++(-0.8,0) coordinate (inv);
  \draw (amp.+) -- ++(-0.8,0) coordinate (noninv);

  % 反馈分压
  \draw (noninv) -- ++(-0.5,0) node[left]{$V_{ref}$};
  \draw (inv) -- ++(-0.5,0) -- ++(0,-2) to[R, l=$R_1$] ++(0,-1.5) node[ground]{};
  \draw (inv) -- ++(-0.5,0) -- ++(0,-2) -- ++(1,0) coordinate (fb);

  % 调整管 PMOS
  \draw (amp.out) -- ++(0.5,0) -- ++(0,1.5) node[pmos, anchor=gate] (pass) {$M_{pass}$};
  \draw (pass.source) -- ++(0,0.8) node[vcc]{$V_{in}$};
  \draw (pass.drain) -- ++(0,-0.5) coordinate (vout);

  % 输出
  \draw (vout) -- ++(1.5,0) node[right]{$V_{out}$};
  \draw (vout) -- ++(0,-0.5) to[C, l=$C_{out}$] ++(0,-1.5) node[ground]{};
  \draw (vout) -- ++(0,-0.5) -- ++(0.8,0) to[R, l=$R_L$] ++(0,-1.5) node[ground]{};

  % 反馈连接
  \draw (vout) -- ++(0,-3) to[R, l=$R_2$] ++(0,-1.5) node[ground]{};
  \draw (vout) -- ++(0,-3) -- ++(-4.3,0) -- (fb);
\end{circuitikz}

\newpage

% ============================================================
% 示例5：Cascode 结构
% ============================================================
\subsection*{5. Cascode 结构}

\begin{circuitikz}[american, scale=1.0, transform shape]
  % M2 (cascode)
  \draw (0,0) node[nmos] (m2) {$M_2$};
  \draw (m2.gate) -- ++(-1,0) node[left]{$V_{bias}$};
  \draw (m2.drain) -- ++(0,0.8) to[R, l=$R_D$] ++(0,1.2) node[vcc]{$V_{DD}$};

  % M1 (输入管)
  \draw (0,-2.5) node[nmos] (m1) {$M_1$};
  \draw (m1.gate) -- ++(-1,0) node[left]{$V_{in}$};
  \draw (m1.source) -- ++(0,-0.5) node[ground]{};

  % 连接 M2 source 到 M1 drain
  \draw (m2.source) -- (m1.drain);

  % 输出
  \draw (m2.drain) -- ++(1.2,0) node[right]{$V_{out}$};
\end{circuitikz}

\vspace{1cm}

% ============================================================
% 示例6：Bandgap 核心（简化）
% ============================================================
\subsection*{6. Bandgap 基准源核心电路}

\begin{circuitikz}[american, scale=0.85, transform shape]
  % 运放
  \draw (0,0) node[op amp] (amp) {};

  % 左侧：R + Q1 (n=1)
  \draw (amp.+) -- ++(-0.8,0) coordinate (pos);
  \draw (pos) -- ++(0,1.5) to[R, l=$R_1$] ++(0,1.5) coordinate (top);
  \draw (top) -- ++(0,0) node[ground, rotate=180, yscale=-1] (q1g) {};
  % 用 npn 代替
  \draw (pos) -- ++(0,1.5) -- ++(-1.5,0) node[npn, anchor=base] (q1) {$Q_1$};
  \draw (q1.collector) -- ++(0,0.5) coordinate (q1c);
  \draw (q1.emitter) -- ++(0,-0.5) node[ground]{};

  % 右侧：R + Q2 (n=8)
  \draw (amp.-) -- ++(-0.8,0) coordinate (neg);
  \draw (neg) -- ++(0,1.5) to[R, l=$R_2$] ++(0,1.5);
  \draw (neg) -- ++(0,1.5) -- ++(1.5,0) node[npn, anchor=base] (q2) {$Q_2$};
  \draw (q2.collector) -- ++(0,0.5) coordinate (q2c);
  \draw (q2.emitter) -- ++(0,-0.5) node[ground]{};

  % 输出到 PMOS 电流镜
  \draw (amp.out) -- ++(1,0) -- ++(0,2) node[pmos, anchor=gate] (p1) {$M_1$};
  \draw (p1.source) -- ++(0,0.8) node[vcc]{$V_{DD}$};
  \draw (p1.drain) -- (q1c);

  \draw (amp.out) -- ++(1,0) -- ++(2,2) node[pmos, anchor=gate] (p2) {$M_2$};
  \draw (p2.source) -- ++(0,0.8) node[vcc]{$V_{DD}$};
  \draw (p2.drain) -- (q2c);

  % 输出
  \draw (p2.drain) -- ++(1.5,0) node[right]{$V_{ref}$};
\end{circuitikz}

\vspace{1cm}

% ============================================================
% 元件速查表
% ============================================================
\section*{Circuitikz 元件速查}

\begin{tabular}{lll}
\hline
元件 & 代码 & 说明 \\
\hline
NMOS 管 & \texttt{node[nmos]} & N沟道MOS \\
PMOS 管 & \texttt{node[pmos]} & P沟道MOS \\
NPN 三极管 & \texttt{node[npn]} & 双极型 \\
电阻 & \texttt{to[R, l=...]} & \\
电容 & \texttt{to[C, l=...]} & \\
电感 & \texttt{to[L, l=...]} & \\
电压源 & \texttt{node[vcc]} & \\
接地 & \texttt{node[ground]} & \\
运放 & \texttt{node[op amp]} & \\
电流源 & \texttt{to[I, l=...]} & \\
\hline
\end{tabular}

\end{document}
